\documentclass[12pt]{article}

\input{../../_assets/preambulo.tex}


\begin{document}

    % 1. Foto de fondo
    % 2. Título
    % 3. Encabezado Izquierdo
    % 4. Color de fondo
    % 5. Coord x del titulo
    % 6. Coord y del titulo
    % 7. Fecha

    
    \input{../../_assets/portada}
    \portadaExamen{ffccA4.jpg}{Modelos\\Matemáticos I\\Examen IV}{MM I. Examen IV}{MidnightBlue}{-8}{28}{2026}{Roxana Acedo Parra}

    \begin{description}
        \item[Asignatura] Modelos Matemáticos I.
        \item[Curso Académico] 2025-26.
        \item[Grado] Doble Grado en Ingeniería Informática y Matemáticas.
        \item[Grupo] Único.
        \item[Profesor] María José Cáceres Granados.
        \item[Descripción] Prueba 1. Temas 0 y 1.
        \item[Fecha] 7 de abril de 2026.
        \item[Duración] 105 minutos.
    
    \end{description}
    \newpage
    
    \begin{ejercicio}
        Un equipo de farmacología trabaja con un modelo matemático propuesto por Ana, la matemática del grupo, para determinar la eficacia de un cierto medicamento frente a una enfermedad bacteriana. Ana les ha dicho que:
        
        \begin{itemize}
        \item Su modelo se basa en una ecuación en diferencias de la forma $p_{n+1} = f(p_n)$, con $f$ una función lo regular que haga falta. Sin entrar en detalles, porque ella es la única que sabe matemáticas.
        
        \item Si la población bacteriana, sujeta a la influencia del medicamento que están estudiando, se rige por este modelo, su tamaño en el día de estudio $n$ viene dado por
            $$ p_n = \frac{1}{a + bn}, \quad a > 0, b > 0, $$
            medido en miles.
        \end{itemize}
        
        \begin{enumerate}
            \item \emph{(1 punto)} ¿Puede ser el modelo planteado una ecuación lineal de la forma $p_{n+1} = Ap_n + B_n$? Razona tu respuesta.
            \item \emph{(1 punto)} Determina el tamaño de la población en función de su tamaño inicial.
            \item \emph{(1 punto)} ¿Es eficaz el medicamento según este modelo? Razona tu respuesta.
            \item \emph{(1 punto)} Según $p_0$, ¿cuántos días son necesarios, como mínimo, para reducir la población a la mitad de la que se tenía en el día 0? Razona tu respuesta.
        \end{enumerate}
    \end{ejercicio}

    \begin{ejercicio}
        Consideramos la ecuación en diferencias:
        \begin{equation}\label{eq:1}
        x_{n+1} = 3x_n + b_n,
        \end{equation}
        donde $\{b_n\}_{n \ge 0}$ es una sucesión de números reales.
    
        \begin{enumerate}
            \item \emph{(1 punto)} Encuentra una expresión, lo más simplificada posible, del término general de las soluciones de la ecuación \eqref{eq:1}.
            \item \emph{(1 punto)} Si $b_n = e^{-n}$, demuestra que $\{y_n\}_{n \ge 0}$, con $y_n = \displaystyle \frac{e^{-n+1}}{1-3e}$, es una solución de~\eqref{eq:1}.
            \item \emph{(1 punto)} Resuelve la ecuación si $b_n = e^{-n}$. ¿Qué ocurre con sus soluciones a largo plazo?
        \end{enumerate}
    \end{ejercicio}

    \begin{ejercicio}
        Se considera la ecuación en diferencias $x_{n+1} = f(x_n)$, donde $f \in C^1(\mathbb{R})$ satisface
        \begin{itemize}
            \item $f(0) = 0$,
            \item $f'(x) > 1 \ \forall x \in \mathbb{R}$.
        \end{itemize}
        
        \begin{enumerate}
            \item \emph{(1 punto)} Calcula los puntos de equilibrio del modelo.
            \item \emph{(1 punto)} Discute la existencia de 2-ciclos no triviales.
            \item \emph{(1 punto)} Si $f(x) = x + e^x - 1$, determina la estabilidad del punto de equilibrio $x = 0$.
        \end{enumerate}
    \end{ejercicio}

    \newpage
	\setcounter{ejercicio}{0}

    \begin{ejercicio}
        Un equipo de farmacología trabaja con un modelo matemático propuesto por Ana, la matemática del grupo, para determinar la eficacia de un cierto medicamento frente a una enfermedad bacteriana. Ana les ha dicho que:
        
        \begin{itemize}
        \item Su modelo se basa en una ecuación en diferencias de la forma $p_{n+1} = f(p_n)$, con $f$ una función lo regular que haga falta. Sin entrar en detalles, porque ella es la única que sabe matemáticas.
        
        \item Si la población bacteriana, sujeta a la influencia del medicamento que están estudiando, se rige por este modelo, su tamaño en el día de estudio $n$ viene dado por
            $$ p_n = \frac{1}{a + bn}, \quad a > 0, b > 0, $$
        medido en miles.
        \end{itemize}
        
        \begin{enumerate}
            \item \emph{(1 punto)} ¿Puede ser el modelo planteado una ecuación lineal de la forma $p_{n+1} = Ap_n + B_n$? Razona tu respuesta. \\
            
            Analicemos la naturaleza del modelo:
                \begin{itemize}
                    \item El enunciado especifica que el modelo sigue una ecuación de la forma $p_{n+1} = f(p_n)$. Esto define un modelo autónomo, donde la función de $f$ depende únicamente del valor de la población en el paso anterior ($p_n$) y no de $n$.
                    \item En una ecuación lineal $p_{n+1} = Ap_n + B_n$, para que el modelo sea autónomo, los coeficientes $A$ y $B_n$ deben ser constantes (es decir, $B_n = B$).
                \end{itemize}
                
                Veamos si $B_n = B$ cte. Para ello $p_{n+1} = Ap_n + B_n$ debería ser compatible con la solución general dada:
                $$ p_n = \frac{1}{a + bn} $$
                $$\frac{1}{a + b(n+1)} = \frac{A}{a + bn} + Bn \implies B_n = \frac{1}{a + b(n+1)} - \frac{A}{a + bn} = \frac{a + bn - A(a + bn + b)}{(a + bn + b)(a + bn)}$$
                Si $A=1$ podríamos deshacernos de la $n$ del numerador, sin embargo el valor del denominador seguirá dependiendo de $n$.
                
                Por tanto, el modelo no puede ser una ecuación lineal autónoma.
            
            \item \emph{(1 punto)} Determina el tamaño de la población en función de su tamaño inicial. \\
                Sea $p_0$ el tamaño de la población inicial, sabemos que:
                 $$p_0 = p_{n=0} = \frac{1}{a+b \cdot 0} \implies p_0 = \frac{1}{a} \implies a = \frac{1}{p_0} $$
                De esta manera la ecucación queda:
                $$p_n = \frac{1}{\frac{1}{p_0} + bn}$$
                
            \item \emph{(1 punto)} ¿Es eficaz el medicamento según este modelo? Razona tu respuesta. \\
                El medicamento será eficaz si a largo plazo elimina la población bacteriana, esto es $\lim\limits_{n \to \infty} p_n = 0$. Veamos si es cierto:
                $$\lim_{n \to \infty} p_n = \lim_{n \to \infty} \frac{1}{\frac{1}{p_0} + bn} = \frac{1}{\infty} = 0$$ 
                Consecuentemente, el medicamento es eficaz según este modelo.

            \item \emph{(1 punto)} Según $p_0$, ¿cuántos días son necesarios, como mínimo, para reducir la población a la mitad de la que se tenía en el día 0? Razona tu respuesta. \\
                Calculamos $n$ para $\nicefrac{p_0}{2}$:
                $$\frac{p_0}{2} = \frac{1}{\frac{1}{p_0} + b \cdot n} \implies \left( \frac{1}{p_0} + b \cdot n \right) \cdot \frac{p_0}{2} = 1 \implies \frac{1}{2} + \frac{p_0 \cdot b \cdot n}{2} = 1$$ 
                $$\implies \frac{p_0 \cdot b \cdot n}{2} = \frac{1}{2} \implies n = \left\lceil \frac{1}{p_0 \cdot b} \right\rceil$$ 
                Se necesitarán $\displaystyle \left\lceil \frac{1}{p_0 \cdot b} \right\rceil$ días como mínimo.
        \end{enumerate}
    \end{ejercicio}

    \begin{ejercicio}
        Consideramos la ecuación en diferencias:
        \begin{equation}\label{eq:2}
        x_{n+1} = 3x_n + b_n,
        \end{equation}
        donde $\{b_n\}_{n \ge 0}$ es una sucesión de números reales.
    
        \begin{enumerate}
            \item \label{item:1}
            \emph{(1 punto)} Encuentra una expresión, lo más simplificada posible, del término general de las soluciones de la ecuación \eqref{eq:2}. \\
            Partiendo de la ecuación en diferencias $x_{n+1} = 3x_n + b_n$, iteramos los primeros términos para identificar un patrón:
            \begin{align*}
                x_1 &= 3x_0 + b_0 \\
                x_2 &= 3x_1 + b_1 = 3(3x_0 + b_0) + b_1 = 3^2x_0 + 3b_0 + b_1 \\
                x_3 &= 3x_2 + b_2 = 3(3^2x_0 + 3b_0 + b_1) + b_2 = 3^3x_0 + 3^2b_0 + 3b_1 + b_2 
            \end{align*}
            Observando la estructura de los términos, proponemos como hipótesis del término general:
            $$x_n = 3^n x_0 + \sum_{i=0}^{n-1} 3^{n-1-i} b_i$$ 
            
            Demostraremos la expresión por inducción sobre $n$:
            \begin{itemize}
                \item Caso base ($n=1$):
                Sustituyendo en la fórmula propuesta:
                $$x_1 = 3^1 x_0 + \sum_{i=0}^{0} 3^{1-1-i} b_i = 3x_0 + 3^0 b_0 = 3x_0 + b_0$$
                Lo cual coincide con la definición de la relación de recurrencia.
                
                \item Paso inductivo:
                Supongamos que la fórmula es válida para $n$ (Hipótesis de Inducción). Veamos si se cumple para $n+1$:
                $$x_{n+1} = 3x_n + b_n$$
                Aplicando la hipótesis de inducción en $x_n$:
                $$ x_{n+1} = 3 \left( 3^n x_0 + \sum_{i=0}^{n-1} 3^{n-1-i} b_i \right) + b_n $$
                Distribuyendo el factor 3:
                $$ x_{n+1} = 3^{n+1} x_0 + \sum_{i=0}^{n-1} 3^{n-i} b_i + b_n $$
                Nótese que el término $b_n$ puede reescribirse como $3^{n-n} b_n$ y añadirlo al sumatorio como el término correspondiente a $i=n$:
                $$ x_{n+1} = 3^{n+1} x_0 + \sum_{i=0}^{n} 3^{n-i} b_i $$
        \end{itemize}
        La fórmula queda demostrada para todo $n \ge 1$.
            \item \label{item:2}
            \emph{(1 punto)} Si $b_n = e^{-n}$, demuestra que $\{y_n\}_{n \ge 0}$, con $y_n = \displaystyle \frac{e^{-n+1}}{1-3e}$, es una solución de \eqref{eq:2}. \\
                Para que $y_n$ sea solución, deberá cumplir \eqref{eq:2}. Entonces, sustituyendo en \eqref{eq:2}:
                $$\frac{e^{-n}}{1-3e} = 3 \frac{e^{-n+1}}{1-3e} + e^{-n} \implies \frac{e^{-n}}{1-3e} = \frac{3 \cdot e \cdot e^{-n}}{1-3e} + e^{-n}$$
                Dividiendo por $e^{-n}$:
                $$\frac{1}{1-3e} = \frac{3e}{1-3e} + 1 = \frac{3e + 1 - 3e}{1-3e} = \frac{1}{1-3e}$$
                De esta manera, $y_n = \displaystyle \frac{e^{-n+1}}{1-3e}$ es solución de \eqref{eq:2}.
            \item\label{item:3}
            \emph{(1 punto)} Resuelve la ecuación si $b_n = e^{-n}$. ¿Qué ocurre con sus soluciones a largo plazo? \\
                Hemos demostrado en el apartado~\ref{item:2} que $y_n = \displaystyle \frac{e^{-n+1}}{1-3e}$ es solución particular de \eqref{eq:2} y en el apartado~\ref{item:1} que las soluciones de \eqref{eq:2} son de la forma $x_n = 3^n(x_0 - \bar{x}_0) + \bar{x}_n$. Consecuentemente:
                $$ x_n = 3^n \left( x_0 - \frac{e}{1-3e} \right) + \frac{e^{-n+1}}{1-3e} $$
                Ahora, $\lim\limits_{n \to \infty} x_n$:
                $$ L = \lim_{n \to \infty} \left[ 3^n \left( x_0 - \frac{e}{1-3e} \right) + \frac{e^{-n+1}}{1-3e} \right] $$
                Como $\lim\limits_{n \to \infty} \displaystyle \frac{e^{-n+1}}{1-3e} = 0$:
                \begin{itemize}
                    \item Si $x_0 = \displaystyle \frac{e}{1-3e}$, sus soluciones convergen al 0 a largo plazo.
                    \item Si $x_0 \neq \displaystyle \frac{e}{1-3e}$, sus soluciones divergen:
                    \begin{itemize}
                        \item Si $x_0 > \displaystyle \frac{e}{1-3e} \implies L = +\infty$ (divergen a $+\infty$).
                        \item Si $x_0 < \displaystyle \frac{e}{1-3e} \implies L = -\infty$ (divergen a $-\infty$).
                    \end{itemize}
                \end{itemize}
        \end{enumerate}
    \end{ejercicio}

    \begin{ejercicio}
        Se considera la ecuación en diferencias $x_{n+1} = f(x_n)$, donde $f \in C^1(\mathbb{R})$ satisface
        \begin{itemize}
            \item $f(0) = 0$,
            \item $f'(x) > 1 \ \forall x \in \mathbb{R}$.
        \end{itemize}
        
        \begin{enumerate}
            \item \emph{(1 punto)} Calcula los puntos de equilibrio del modelo. \\
                Los puntos de equilibrio son las soluciones de la ecuación $f(x) = x$. Por el enunciado, sabemos que $f(0) = 0$, entonces $x = 0$ es un punto de equilibrio. 
    
                Para determinar si existen otros puntos fijos, definimos la función auxiliar $g(x) = f(x) - x$. Es evidente que $g(0) = 0$. Derivando respecto a $x$, obtenemos:
                $$g'(x) = f'(x) - 1$$
                Dado que $f'(x) > 1$ para todo $x \in \mathbb{R}$, se sigue que $g'(x) > 0$. Esto implica que $g(x)$ es una función \textbf{estrictamente creciente} en todo su dominio.
                
                Supongamos ahora que existe otro punto de equilibrio $y \neq 0$, tal que $g(y) = 0$. Analizamos los dos casos posibles:
                \begin{itemize}
                    \item Si $y < 0$, por el crecimiento estricto de $g$, tendríamos $g(y) < g(0)$, lo que implica $0 < 0$, una contradicción.
                    \item Si $y > 0$, análogamente tendríamos $g(y) > g(0)$, lo que implica $0 > 0$, de nuevo una contradicción.
                \end{itemize}
                
                Por lo tanto, $g(x)$ se anula únicamente en $x = 0$, lo que demuestra que $\exists!$ punto de equilibrio en el modelo y es $x = 0$.
                
            \item \emph{(1 punto)} Discute la existencia de 2-ciclos no triviales. \\
                Supongamos, por reducción al absurdo, que existe un 2-ciclo no trivial $\{x_0, x_1\}$, lo que implica que $f(x_0) = x_1$ y $f(x_1) = x_0$ con $x_0 \neq x_1$. Sin pérdida de generalidad, podemos asumir que $x_0 < x_1$.
    
                Dado que el enunciado establece que $f'(x) > 1$ para todo $x \in \mathbb{R}$, la función $f$ es estrictamente creciente. Aplicando la propiedad de monotonía a nuestra hipótesis:
                $$x_0 < x_1 \implies f(x_0) < f(x_1)$$
                Sustituyendo los valores del ciclo en la desigualdad obtenemos:
                $$x_1 < x_0$$
                Esto representa una contradicción con la suposición inicial $x_0 < x_1$. Por lo tanto, concluimos que no pueden existir 2-ciclos no triviales bajo estas condiciones.
                
            \item \emph{(1 punto)} Si $f(x) = x + e^x - 1$, determina la estabilidad del punto de equilibrio $x = 0$. \\
                En primer lugar, verificamos que $x = 0$ es, en efecto, un punto fijo evaluando la función:
                $$f(0) = 0 + e^0 - 1 = 0 + 1 - 1 = 0$$
                Para determinar su estabilidad, calculamos la derivada de la función:
                $$f'(x) = 1 + e^x$$
                Evaluando en el punto de equilibrio obtenemos $f'(0) = 1 + e^0 = 2$. Nótese que la condición del enunciado se cumple, ya que $f'(x) > 1$ para todo $x \in \mathbb{R}$ dado que $e^x > 0$.
                
                Finalmente, aplicando el teorema de la primera derivada, tenemos que:
                $$|f'(0)| = 2 > 1$$
                entonces el punto de equilibrio $x = 0$ es inestable.

        \end{enumerate}
    \end{ejercicio}

\end{document}
